\documentclass[12pt,a4paper]{article}

\usepackage[utf8]{inputenc}
\usepackage[T1]{fontenc}
\usepackage{lmodern}
\usepackage{textcomp}
\usepackage{newunicodechar}
\usepackage[margin=2.5cm]{geometry}
\usepackage{amsmath,amssymb}
\usepackage{graphicx}
\usepackage{listings}
\usepackage{xcolor}
\usepackage{hyperref}
\usepackage{fancyhdr}
\usepackage{float}
\usepackage{booktabs}
\usepackage{caption}
\usepackage{tabularx}
\usepackage{array}

% Adjust header height for fancyhdr
\setlength{\headheight}{15pt}
\addtolength{\topmargin}{-3pt}

% Romanian Unicode support for standard pdflatex / tectonic
\newunicodechar{Ă}{\v{A}}
\newunicodechar{ă}{\v{a}}
\newunicodechar{Â}{\^{A}}
\newunicodechar{â}{\^{a}}
\newunicodechar{Î}{\^{I}}
\newunicodechar{î}{\^{i}}
\newunicodechar{Ș}{\c{S}}
\newunicodechar{ș}{\c{s}}
\newunicodechar{Ț}{\c{T}}
\newunicodechar{ț}{\c{t}}

\hypersetup{
    colorlinks=true,
    linkcolor=blue,
    filecolor=magenta,
    urlcolor=cyan,
    citecolor=blue
}

% Colors for code listings
\definecolor{codebg}{gray}{0.96}
\definecolor{codegreen}{rgb}{0.13,0.55,0.13}
\definecolor{codegray}{rgb}{0.5,0.5,0.5}
\definecolor{codepurple}{rgb}{0.58,0,0.82}
\definecolor{codeblue}{rgb}{0.05,0.25,0.70}

\lstset{
    backgroundcolor=\color{codebg},
    basicstyle=\ttfamily\footnotesize,
    breaklines=true,
    breakatwhitespace=true,
    frame=single,
    rulecolor=\color{gray!40},
    keepspaces=true,
    showstringspaces=false,
    upquote=true,
    columns=fullflexible,
    tabsize=2
}

\lstdefinestyle{cstyle}{
    language=C,
    commentstyle=\color{codegreen}\itshape,
    keywordstyle=\color{codeblue}\bfseries,
    stringstyle=\color{codepurple},
    numbers=left,
    numbersep=7pt,
    numberstyle=\tiny\color{codegray},
    stepnumber=1
}

\lstdefinestyle{bashstyle}{
    language=bash,
    commentstyle=\color{codegreen}\itshape,
    keywordstyle=\color{codeblue}\bfseries,
    stringstyle=\color{codepurple},
    numbers=left,
    numbersep=7pt,
    numberstyle=\tiny\color{codegray},
    stepnumber=1
}

\lstdefinestyle{consolestyle}{
    backgroundcolor=\color{codebg},
    basicstyle=\ttfamily\footnotesize,
    breaklines=true,
    frame=single,
    rulecolor=\color{gray!40},
    numbers=none
}

% Header and Footer
\pagestyle{fancy}
\fancyhf{}
\lhead{Operating Systems -- Lab 13}
\rhead{Iulian Pavlov, FAF-241}
\cfoot{\thepage}

\begin{document}

% ============================================================================
% TITLE PAGE
% ============================================================================
\hypersetup{pageanchor=false}
\begin{titlepage}
\begin{center}
    \includegraphics[width=0.3\textwidth]{logo.png} \\[0.2cm]

    \textbf{MINISTRY OF EDUCATION, CULTURE AND RESEARCH \\ OF THE REPUBLIC OF MOLDOVA} \\[0.3cm]

    \textbf{Technical University of Moldova \\ Faculty of Computers, Informatics and Microelectronics \\ Department of Software and Automation Engineering} \\[1.8cm]

    \textbf{Iulian Pavlov, FAF-241}\\[0.5cm]

    \Huge \textbf{Report} \\[0.5cm]

    \large Laboratory work nr. 13 \\[0.5cm]

    \textbf{\Large Semaphores in Your Own Kernel} \\[0.3cm]
    \normalsize Course: Operating Systems (Track B: build your own small OS (xv6)) \\[2.2cm]

    \begin{flushright}
        \textit{Checked by:} \\
        \textbf{Reițman Patricia}, \textit{University Assistant} \\
        Department of Software and Automation Engineering, \\
        FCIM Faculty, UTM
    \end{flushright}
    \vfill

    Chișinău -- 2026
\end{center}
\end{titlepage}
\hypersetup{pageanchor=true}

\pagenumbering{roman}
\newpage
\tableofcontents
\newpage
\pagenumbering{arabic}

% ============================================================================
% 1. THEORETICAL BACKGROUND & OBJECTIVES
% ============================================================================
\section{Theoretical Background and Objectives}

\subsection{Laboratory Objectives and Scope}
Whereas Laboratories No.~11 and~12 exercised POSIX semaphores from user space on Linux, Laboratory Work No.~13 moves inside the \textbf{xv6-riscv} kernel to build counting semaphores from scratch. Each semaphore is simply an integer counter protected by a per-semaphore kernel spinlock and paired with sleep/wakeup process scheduling so blocked waiters yield the CPU cleanly.

The objectives of Laboratory Work No.~13 are:
\begin{itemize}
    \item \textbf{User-Space Syscall Plumbing}: Assigning syscall numbers \texttt{24..26} (\texttt{SYS\_\allowbreak sem\_\allowbreak init}, \texttt{SYS\_\allowbreak sem\_\allowbreak down}, \texttt{SYS\_\allowbreak sem\_\allowbreak up}) alongside \texttt{SYS\_\allowbreak trace} (\texttt{23}), adding prototypes in \texttt{user/\allowbreak user.h}, emitting RISC-V \texttt{ecall} stubs in \texttt{user/\allowbreak usys.pl}, adding \texttt{semtest.c} and \texttt{mix.c}, and verifying Check~1.
    \item \textbf{Kernel Semaphore Table (\texttt{kernel/\allowbreak sem.c})}: Implementing \texttt{NSEM = 8} semaphores protected by spinlocks and the two-step calls \texttt{sleep\_\allowbreak prepare(s)}, \texttt{release(\allowbreak\&s->\allowbreak lock)}, and \texttt{sleep()}.
    \item \textbf{Verification and Mutual Exclusion}: Passing both ordering and ping-pong tests in \texttt{semtest}, contrasting unlocked vs.\ semaphore-locked output in \texttt{mix} across 3 QEMU CPUs, implementing optional \texttt{sem\_\allowbreak trydown} (syscall \texttt{27}), and committing the working tree to git.
\end{itemize}

\subsection{Two-Step Sleep Protocol in Our xv6 Kernel}
Unlike textbook xv6, which passes the condition lock directly to \texttt{sleep(chan, lk)}, this xv6 tree (\texttt{kernel/\allowbreak proc.c}) separates channel registration from parking the thread:
\begin{enumerate}
    \item \texttt{sleep\_\allowbreak prepare(s)} acquires \texttt{p->lock}, sets \texttt{p->chan = s}, and releases \texttt{p->lock} \emph{while \texttt{s->lock} is still held}, so no \texttt{wakeup(s)} can slip past unseen before registration.
    \item \texttt{release(\&s->lock)} releases the semaphore spinlock so another hart can call \texttt{sem\_\allowbreak up()} and \texttt{mycpu()->noff} drops from \texttt{1} back to \texttt{0}.
    \item \texttt{sleep()} re-acquires \texttt{p->lock} and checks \texttt{if (p->chan != 0)}. If \texttt{wakeup(s)} arrived in the window between steps~1 and~3, \texttt{wakeup()} already cleared \texttt{p->chan = 0}, so \texttt{sleep()} returns immediately without blocking; otherwise \texttt{sleep()} sets \texttt{p->state = SLEEPING} and calls \texttt{sched()}.
\end{enumerate}

Table~\ref{tab:syscalls} summarizes all new semaphore system calls registered in \texttt{kernel/\allowbreak syscall.h}.

\begin{table}[H]
\centering
\caption{Semaphore System Calls Implemented in Lab 13}
\label{tab:syscalls}
\begin{tabularx}{\textwidth}{l c l >{\raggedright\arraybackslash}X}
\toprule
\textbf{Constant} & \textbf{Number} & \textbf{User Prototype} & \textbf{Kernel Behavior} \\
\midrule
\texttt{SYS\_\allowbreak sem\_\allowbreak init} & 24 & \texttt{int sem\_\allowbreak init(int, int);} & Initializes \texttt{sems[id].value = val} (if valid) \\
\texttt{SYS\_\allowbreak sem\_\allowbreak down} & 25 & \texttt{int sem\_\allowbreak down(int);} & Decrements if \texttt{> 0}, or sleeps via \texttt{sleep\_\allowbreak prepare/sleep} \\
\texttt{SYS\_\allowbreak sem\_\allowbreak up} & 26 & \texttt{int sem\_\allowbreak up(int);} & Increments \texttt{value} and calls \texttt{wakeup(s)} \\
\texttt{SYS\_\allowbreak sem\_\allowbreak trydown} & 27 & \texttt{int sem\_\allowbreak trydown(int);} & Optional: decrements if \texttt{> 0}, else returns \texttt{-1} at once \\
\bottomrule
\end{tabularx}
\end{table}

% ============================================================================
% 2. PART 1: USER-SIDE PLUMBING AND CHECK 1
% ============================================================================
\section{Part 1: User-Side Plumbing and Test Programs}

\subsection{Syscall Numbers, Prototypes, Assembly Stubs, and Makefile}
Keeping \texttt{SYS\_\allowbreak sync = 22} and \texttt{SYS\_\allowbreak trace = 23} intact from Lab~10, the user-space plumbing adds:
\begin{itemize}
    \item \textbf{\texttt{kernel/\allowbreak syscall.h}}: Defines \texttt{SYS\_\allowbreak sem\_\allowbreak init} (\texttt{24}), \texttt{SYS\_\allowbreak sem\_\allowbreak down} (\texttt{25}), \texttt{SYS\_\allowbreak sem\_\allowbreak up} (\texttt{26}), and optional \texttt{SYS\_\allowbreak sem\_\allowbreak trydown} (\texttt{27}).
    \item \textbf{\texttt{user/\allowbreak user.h} \& \texttt{user/\allowbreak usys.pl}}: Declares the C wrappers and adds \texttt{entry("sem\_\allowbreak init")}, \texttt{entry("sem\_\allowbreak down")}, \texttt{entry("sem\_\allowbreak up")}, and \texttt{entry("sem\_\allowbreak trydown")} so \texttt{perl user/\allowbreak usys.pl} emits RISC-V \texttt{ecall} stubs into \texttt{user/\allowbreak usys.S}.
    \item \textbf{\texttt{Makefile} (\texttt{UPROGS})}: Adds \texttt{\$U/\_semtest\textbackslash}, \texttt{\$U/\_mix\textbackslash}, and optional \texttt{\$U/\_trydown\textbackslash} directly after \texttt{\$U/\_zombie\textbackslash}.
\end{itemize}

\subsection{Test Programs: \texttt{user/semtest.c} and \texttt{user/mix.c}}
\lstinputlisting[style=cstyle, caption={Ordering and ping-pong test suite (\texttt{user/semtest.c})}]{semtest.c}

\newpage
\lstinputlisting[style=cstyle, caption={Concurrent console write test without and with semaphore 3 (\texttt{user/mix.c})}]{mix.c}

\newpage
\subsection{Check 1: Behavior Before Kernel-Side Semaphores}
Building and booting QEMU after Part~1 (before wiring \texttt{kernel/\allowbreak sem.c}) confirmed that the stubs reached kernel trap dispatch:
\begin{lstlisting}[style=consolestyle, caption={Observed Check 1 output before implementing kernel semaphores}]
$ semtest
3 semtest: unknown sys call 24
4 semtest: unknown sys call 25
child: I ran after the parent
parent: done, I let the child go
3 semtest: unknown sys call 26
3 semtest: unknown sys call 24
...
pi npongg  11
...
ppoinngg  22
...
 ppoinngg  33
...
semtest: OK
\end{lstlisting}
As expected for Check~1, unknown syscalls \texttt{24} and \texttt{25} returned \texttt{-1} immediately, so the child printed \texttt{child: I ran after the parent} \emph{before} the parent finished \texttt{pause(10)}, and Test~2 interleaved \texttt{pi npongg 11} concurrently.

% ============================================================================
% 3. PART 2: THE KERNEL SIDE
% ============================================================================
\newpage
\section{Part 2: Kernel Semaphore Implementation}

\subsection{Semaphore Table and Handlers (\texttt{kernel/sem.c})}
\texttt{kernel/\allowbreak sem.c} allocates a static table of \texttt{NSEM = 8} \texttt{struct sem} entries shared globally across all processes. Zero-initialization in BSS starts every \texttt{spinlock} unlocked. A small shared helper \texttt{fetch\_\allowbreak sem()} extracts argument~\texttt{0} and bounds-checks \texttt{0 <= id < NSEM} across all four handlers.

\lstinputlisting[style=cstyle, basicstyle=\ttfamily\scriptsize, caption={Kernel counting semaphores with optional \texttt{sem\_trydown} (\texttt{kernel/sem.c})}]{sem.c}

\subsection{Kernel Makefile, Dispatch Table, and Git Commit}
To compile \texttt{kernel/\allowbreak sem.c} into the kernel binary and register all syscall handlers:
\begin{itemize}
    \item \textbf{\texttt{Makefile} (\texttt{OBJS})}: Inserted \texttt{\$K/\allowbreak sem.o \textbackslash} directly after \texttt{\$K/\allowbreak sysproc.o \textbackslash}.
    \item \textbf{\texttt{kernel/\allowbreak syscall.c}}: Declared \texttt{sys\_\allowbreak sem\_\allowbreak *()} prototypes, mapped \texttt{SYS\_\allowbreak sem\_\allowbreak init} (\texttt{24}) through \texttt{SYS\_\allowbreak sem\_\allowbreak trydown} (\texttt{27}) in \texttt{syscalls[]}, and added matching names in \texttt{syscallnames[]} so Lab~10's \texttt{trace} facility never reads out of bounds.
    \item \textbf{Git Commit (\texttt{\textasciitilde/xv6-riscv})}: Committed the completed working tree to branch \texttt{riscv} at commit \texttt{aae203a}.
\end{itemize}

\newpage
\subsection{Check 2: Verified Ordering and Ping-Pong (\texttt{semtest})}
After rebuilding \texttt{kernel/\allowbreak kernel} and \texttt{fs.img} cleanly with \texttt{make}, running \texttt{semtest} inside 3-CPU QEMU produced:
\begin{lstlisting}[style=consolestyle, caption={Observed Check 2 output with kernel semaphores enabled}]
$ semtest
parent: done, I let the child go
child: I ran after the parent
ping 1
 pong 1
ping 2
 pong 2
ping 3
 pong 3
semtest: OK
\end{lstlisting}
Even though the parent sleeps for \texttt{10} ticks via \texttt{pause(10)}, the child genuinely blocks in \texttt{sys\_\allowbreak sem\_\allowbreak down(0)} until the parent calls \texttt{sem\_\allowbreak up(0)}, and \texttt{ping}/\texttt{pong} alternate in lockstep across all three iterations.

% ============================================================================
% 4. PART 3: MIX, OPTIONAL TRYDOWN, AND OBSERVE ANSWERS
% ============================================================================
\section{Part 3: \texttt{mix}, \texttt{sem\_trydown}, and Analysis}

\subsection{Running \texttt{mix} Without and With Semaphore Locking}
Running \texttt{mix} twice without arguments lets parent and child interleave concurrently on 3~CPUs; passing an argument (\texttt{mix lock}) sets \texttt{uselock = 1} and wraps each 31-byte line in \texttt{sem\_\allowbreak down(3)} / \texttt{sem\_\allowbreak up(3)}:

\begin{lstlisting}[style=consolestyle, caption={Observed QEMU outputs of \texttt{mix} without locking (two runs)}]
$ mix
AAAAAAAAAbAbbbbAbbbAbAbAbAbAbAbAbAbAbAbAbAbbAbAbAbAbAAbAb
bAbAbA
AbAbAbAAbAbbbAAbAbAbAbbAbbAbAbAbAbAbAbAbAbAbAbAAAbAbAb
bbAbA
AbAbAbAbbAbAbAbAbAbAbAbAbAbAbAbAbAbAbAbAbAbAbAAbAbAbAbAb
bAbA
AbAAAbAbAbAbAAAAbAAbbbAbbbAbAbAbAbAbAAbAAbAbAAAbAb
bAbAbAbAbAbAbAbA
AbAbAbAbAbAbbbAbAbAbAbAbAbAbAbAbAAbAbAbAbbAbA
bbbbbbb
$ mix
AAAAAbAbAbAbAbAbAbAbAbAbAbAbAbbAbAbAbAbAbAbAbAbAbAbAbbAA
bAAAAAbAbAAAAAAbA
bbbAbbbAAbAbAbbAbAbAbAbAbAbAbAbAbAbAbAb
bAbAbAbAbAbAbAbA
AbAbAbAbAbAbAbAbAbAbAAbAAbbAbAbAbAbAbAbAbAb
bAbAbAbAbAbAbAbbA
AbAAbAbbbAbAbAbAbAbAbAbAbAbAbAbAbAbAbAbAbAbAb
bAbAbAbAbAbAbA
AbAbAbAbAbAbAbAbAbAbAbAbAbAAbbbAbbAAbAbAbAbAbAbbb
bbbb
\end{lstlisting}

\newpage

\begin{lstlisting}[style=consolestyle, caption={Observed QEMU output of \texttt{mix lock} (semaphore 3 initialized to 1)}]
$ mix lock
AAAAAAAAAAAAAAAAAAAAAAAAAAAAAA
AAAAAAAAAAAAAAAAAAAAAAAAAAAAAA
AAAAAAAAAAAAAAAAAAAAAAAAAAAAAA
AAAAAAAAAAAAAAAAAAAAAAAAAAAAAA
AAAAAAAAAAAAAAAAAAAAAAAAAAAAAA
bbbbbbbbbbbbbbbbbbbbbbbbbbbbbb
bbbbbbbbbbbbbbbbbbbbbbbbbbbbbb
bbbbbbbbbbbbbbbbbbbbbbbbbbbbbb
bbbbbbbbbbbbbbbbbbbbbbbbbbbbbb
bbbbbbbbbbbbbbbbbbbbbbbbbbbbbb
\end{lstlisting}

\subsection{Optional Extension: \texttt{sem\_trydown} (Syscall 27)}
As specified in the optional exercise, \texttt{sys\_\allowbreak sem\_\allowbreak trydown()} (\texttt{SYS\_\allowbreak sem\_\allowbreak trydown = 27}) acquires \texttt{s->lock}, returns \texttt{-1} immediately when \texttt{s->value == 0} instead of calling \texttt{sleep()}, and decrements/returns \texttt{0} when a unit is available. A dedicated user program \texttt{user/\allowbreak trydown.c} verifies both outcomes:

\lstinputlisting[style=cstyle, caption={Non-blocking semaphore test program (\texttt{user/trydown.c})}]{trydown.c}

\begin{lstlisting}[style=consolestyle, caption={Observed QEMU output of optional \texttt{trydown}}]
$ trydown
first trydown: 0
second trydown (at 0): -1
after sem_up, third trydown: 0
\end{lstlisting}

\newpage
\subsection{Observe Answers (Saved in \texttt{notes.txt})}
\begin{enumerate}
    \item \textbf{Without the lock, why do the letters mix? (How many system calls does one line cost?)} \\
    Each call to \texttt{lines()} in \texttt{user/\allowbreak mix.c} prints \texttt{30} identical letters plus \texttt{'\textbackslash n'} (\texttt{31} characters total) one character at a time via \texttt{for (char *p = s; *p; p++) write(1, p, 1);}. Therefore, \textbf{one line costs \texttt{31} \texttt{write()} system calls}, not a single atomic write. When \texttt{mix} runs without arguments, QEMU boots xv6 with \texttt{CPUS := 3}, allowing parent and child processes to run truly in parallel across separate CPUs (and interleave across kernel/timer transitions between individual syscalls), mixing their 1-byte console writes unpredictably.

    \item \textbf{With the lock every line is whole. Which two lines of \texttt{mix.c} make one line a critical section?} \\
    Inside \texttt{lines()} (\texttt{user/\allowbreak mix.c}), \texttt{sem\_\allowbreak down(3);} (immediately before the \texttt{for (char *p = s; *p; p++)} loop, under \texttt{if (uselock)}) and \texttt{sem\_\allowbreak up(3);} (immediately after the loop, under \texttt{if (uselock)}) bracket all \texttt{31} character writes of one line into a single mutually exclusive critical section.

    \item \textbf{In \texttt{sem\_down} the test is \texttt{while (s->value == 0)}, not \texttt{if}. Why can a process wake up and still find \texttt{0}?} \\
    Inside \texttt{sys\_\allowbreak sem\_\allowbreak up()}, \texttt{wakeup(s)} scans \texttt{proc[]} (\texttt{kernel/\allowbreak proc.c}) and marks \emph{every} process sleeping on channel \texttt{s} as \texttt{RUNNABLE}, not just one waiter. If multiple waiters wake together, the first woken process to re-acquire \texttt{s->lock} exits the loop and decrements \texttt{s->value} back from \texttt{1} to \texttt{0}, so any second woken process observes \texttt{s->value == 0} when it next acquires \texttt{s->lock}. Likewise, an already-running process on another hart can claim the unit first ("barging"), or \texttt{kkill()} can mark a sleeping process \texttt{RUNNABLE} without changing \texttt{s->value}. Looping with \texttt{while} guarantees \texttt{s->value} never drops below \texttt{0}.

    \item \textbf{\texttt{sem\_down} calls \texttt{release(\&s->lock)} BEFORE \texttt{sleep()}. What would go wrong if it slept holding the spinlock?} \\
    Two catastrophic failures occur:
    \begin{itemize}
        \item \textbf{Immediate kernel panic (\texttt{sched locks}):} \texttt{sleep()} acquires \texttt{p->lock} and calls \texttt{sched()} (\texttt{kernel/\allowbreak proc.c}), which enforces \texttt{if (mycpu()->noff != 1) panic("sched locks");}. Holding both \texttt{s->lock} and \texttt{p->lock} pushes \texttt{noff} to \texttt{2}, triggering an instant kernel panic.
        \item \textbf{Deadlock / spinlock stall:} Spinlocks disable interrupts via \texttt{push\_\allowbreak off()} (\texttt{kernel/\allowbreak spinlock.c}) and must be released before sleeping. Even without the \texttt{panic()} check, any caller of \texttt{sem\_\allowbreak up()} must run \texttt{acquire(\&s->lock)} before incrementing \texttt{s->value} and calling \texttt{wakeup(s)}; if the sleeping process retained \texttt{s->lock}, no \texttt{sem\_\allowbreak up()} caller on any CPU could ever acquire \texttt{s->lock} to wake it up.
    \end{itemize}
\end{enumerate}

% ============================================================================
% 5. CONCLUSIONS
% ============================================================================
\newpage
\section{Conclusions}

Implementing semaphores inside \textbf{xv6-riscv} connects every user-space abstraction practiced in Labs~11 and~12 down to kernel hardware and scheduler primitives:
\begin{itemize}
    \item A semaphore reduces cleanly to three core pieces (\texttt{kernel/\allowbreak sem.c}): a counter \texttt{value}, a protecting spinlock \texttt{lock}, and sleeping/waking via \texttt{sleep\_\allowbreak prepare(s)}, \texttt{sleep()}, and \texttt{wakeup(s)}.
    \item \texttt{semtest} verifies both one-shot parent-before-child ordering and strict multi-round ping-pong alternation (\texttt{semtest: OK}), replacing CPU-wasting busy-waiting with true blocking in \texttt{SLEEPING} state.
    \item \texttt{mix} demonstrates how \texttt{31} separate 1-byte \texttt{write()} calls per line race across 3~CPUs without locking, while wrapping the loop between \texttt{sem\_\allowbreak down(3)} and \texttt{sem\_\allowbreak up(3)} keeps every 30-character line intact.
    \item Re-checking \texttt{while (s->value == 0)} after waking prevents negative counts under thundering-herd wakeups, and releasing \texttt{s->lock} before \texttt{sleep()} preserves xv6's \texttt{mycpu()->noff == 1} scheduler invariant and prevents deadlock against \texttt{sem\_\allowbreak up()}.
\end{itemize}

\end{document}
